% ============================================================
% Results and Discussion
% All numbers verified against inference_results_thesis.xlsx
% (sheets: english / Spanish / Arabic / Invarience / Analysis)
% and independently recomputed from the per-language sheets.
%
% FIGURES: requires \usepackage{graphicx} and \usepackage{svg}.
% Figure files (SVG): results-tradeoff, results-agreement,
% results-language, results-conflicting, results-lambda.
% TABLES: condensed summary here (tables/table-results-summary.tex);
%   full per-model tables live in chapters/appendix.tex.
%   Requires \usepackage{booktabs}.
% ============================================================

\section{Results and Discussion}
\label{sec:results}

This chapter reports and interprets the experimental results. We first describe the evaluation setup
and the data-hygiene decisions that determine which runs enter the analysis
(Section~\ref{subsec:res-setup}). We then present task performance
(Section~\ref{subsec:res-task}) and cross-lingual consistency
(Section~\ref{subsec:res-consistency}) separately, as the methodology requires, before bringing them
together to answer whether optimizing for language invariance improves ranking
(Section~\ref{subsec:res-invariance}, RQ1). Section~\ref{subsec:res-language} examines where the
cross-lingual differences actually live, Section~\ref{subsec:res-ablations} covers the invariance
weight, numeric normalization, and backbone families, and
Section~\ref{subsec:res-relationship} analyses the relationship between invariance and reasoning
quality (RQ2). Section~\ref{subsec:res-threats} states the threats to validity, including a context-budget
confound that qualifies the accuracy comparison, and Section~\ref{subsec:res-summary} summarizes the
findings against the three research questions (RQ3).

\subsection{Experimental Scope and Data Hygiene}
\label{subsec:res-setup}

All results are computed on the held-out \emph{test} set of the shared task, in the aligned trilingual
form described in Section~\ref{subsec:method-data}: $4{,}233$ claims present in English, Spanish, and
Arabic with parallel trace sets, evaluated per language and, for consistency measures, across the
three language-conditioned outputs of the same claim. The test set is disjoint from the training and
validation splits: no configuration was trained on it, and model selection (early stopping, best-epoch
choice) used the validation split only. The numbers below therefore measure generalization to unseen
claims rather than fit. Every reported configuration was trained and
evaluated with three random seeds and we report the mean. In total we evaluate $32$ configurations
spanning prompted baselines, the earlier trace scorers, standard multilingual fine-tuning, and the
invariance-trained models introduced in this thesis.

\paragraph{Excluded runs.}
Four of the $32$ configurations are excluded from every statistic reported below. The rows
\emph{Mistral, English DPO}, \emph{Llama, inv.\ $\lambda^\star$=0.05}, and \emph{Qwen, prompted} carry
scores that are identical to one another to four decimal places and identical across all three
languages, together with cross-lingual agreement of exactly $1.000$ on every measure. Three distinct
systems cannot plausibly produce byte-identical scores; these rows are duplicated or failed runs
rather than results, and reporting them would contaminate every group mean. We exclude them and
report the analysis over the remaining $28$ configurations.

The fourth exclusion is different in kind and worth stating explicitly, because it illustrates the
central measurement principle of this thesis. The \emph{random ranking} baseline also achieves
perfect cross-lingual agreement---Fleiss $\kappa = 1.000$ and Kendall $\tau = 1.000$ on every language
pair. This is not an artefact: the baseline ranks traces in a fixed seeded permutation of trace
indices that does not depend on the input at all, so it produces literally the same ranking in every
language and is \emph{invariant by construction}. It is also useless, sitting at the bottom of every
accuracy measure. A system can therefore be perfectly language-invariant and still carry no
information, which is precisely why invariance must be reported alongside, and never instead of, task
accuracy~\cite{elazar2021measuringimprovingconsistencypretrained,Qi_2023}. We exclude the random
baseline from correlation statistics for the same reason---it would dominate them while being
uninformative---but retain it in Table~\ref{tab:results-summary} as the accuracy floor.

\paragraph{Achievable ceiling.}
As established in Section~\ref{subsec:method-task}, Recall@5 is structurally capped. On this
test set the upper bound is $0.3005$, and $578$ of the $4{,}233$ claims ($13.7\%$) have no
candidate trace whose verdict matches the gold label, contributing zero Recall@5 regardless of the
ranking. All Recall@5 figures below should be read against this ceiling rather than against $1.0$.

\paragraph{Terminology.}
The task is scored on two metrics that behave differently, so we keep them apart throughout. We use
\emph{verdict quality} for the claim-level macro-F1 and \emph{ranking quality} for Recall@5, and where
a single summary of task performance is needed---for example on the horizontal axis of
Figure~\ref{fig:results-tradeoff}, or in the group comparisons---we use \emph{mean macro-F1}, the
verdict macro-F1 averaged over English, Spanish, and Arabic. Where the word \emph{accuracy} appears
below it is shorthand for this quantity, never for Recall@5 and never for the official combined score.
\emph{Consistency} always refers to the cross-lingual agreement measures (Fleiss $\kappa$, Cohen's
$\kappa$, Kendall $\tau$), which are computed independently of both task metrics.

\subsection{Task Performance}
\label{subsec:res-task}

Table~\ref{tab:results-summary} condenses the results: for each training group it reports the group
mean together with that group's most accurate and most cross-lingually stable member, placing accuracy
and consistency side by side. Per-model results for all $28$ configurations are given in
Appendix~\ref{app:results-full} (Tables~\ref{tab:results-main} and~\ref{tab:results-consistency}).
Several observations follow.

\begin{table}[htbp]
\centering\footnotesize\setlength{\tabcolsep}{3.6pt}
\begin{tabular}{l c rrrr rrrr rr}
\toprule
& & \multicolumn{4}{c}{Recall@5 (ceiling $0.3005$)} & \multicolumn{4}{c}{Macro-F1} & \multicolumn{2}{c}{Consistency} \\
\cmidrule(lr){3-6}\cmidrule(lr){7-10}\cmidrule(lr){11-12}
Configuration & $n$ & EN & ES & AR & Mean & EN & ES & AR & Mean & $\kappa$ & $\tau$ \\
\midrule
\itshape Invariance-trained (this thesis) & 13 & \itshape 0.2673 & \itshape 0.2676 & \itshape 0.2617 & \itshape 0.2656 & \itshape 0.6183 & \itshape 0.6176 & \itshape 0.6085 & \itshape 0.6148 & \itshape 0.9443 & \itshape 0.5667 \\
\quad Qwen, inv.\ $\lambda^\star$=0.1 {\scriptsize (acc.)} &  & 0.2663 & 0.2667 & 0.2622 & 0.2651 & 0.6192 & 0.6197 & 0.6157 & 0.6182 & 0.9479 & 0.5820 \\
\quad Llama, inv.\ $\lambda^\star$=0.3 {\scriptsize (stab.)} &  & 0.2683 & 0.2684 & 0.2638 & 0.2668 & 0.6165 & 0.6150 & 0.6034 & 0.6116 & 0.9428 & \textbf{0.6440} \\
\addlinespace[3pt]
\itshape Standard multilingual fine-tuning & 9 & \itshape 0.2657 & \itshape 0.2728 & \itshape 0.2653 & \itshape 0.2679 & \itshape 0.6158 & \itshape 0.6224 & \itshape 0.5847 & \itshape 0.6076 & \itshape 0.9028 & \itshape 0.1308 \\
\quad Llama, standard {\scriptsize (acc.)} &  & 0.2649 & 0.2704 & \textbf{0.2692} & 0.2682 & 0.6256 & \textbf{0.6320} & 0.6105 & \textbf{0.6227} & 0.8810 & 0.0863 \\
\quad Qwen, num.\,norm. {\scriptsize (stab.)} &  & 0.2702 & 0.2762 & 0.2654 & 0.2706 & 0.6221 & 0.6258 & 0.5810 & 0.6096 & 0.9272 & 0.2496 \\
\addlinespace[3pt]
\itshape Earlier trace scorers & 4 & \itshape 0.2271 & \itshape 0.2406 & \itshape 0.2276 & \itshape 0.2318 & \itshape 0.5861 & \itshape 0.5878 & \itshape 0.5851 & \itshape 0.5863 & \itshape 0.8596 & \itshape 0.0578 \\
\quad Mistral, multiling.\ trace scorer (10 ep.) {\scriptsize (acc.)} &  & 0.2566 & 0.2549 & 0.2520 & 0.2545 & 0.5967 & 0.5932 & 0.5828 & 0.5909 & 0.9040 & 0.0955 \\
\quad Mistral, numeric trace scorer {\scriptsize (stab.)} &  & 0.2184 & 0.2155 & 0.2228 & 0.2189 & 0.5826 & 0.5693 & 0.5938 & 0.5819 & 0.8643 & 0.1075 \\
\addlinespace[3pt]
\itshape Prompted baselines & 2 & \itshape 0.2229 & \itshape 0.2198 & \itshape 0.2174 & \itshape 0.2200 & \itshape 0.5761 & \itshape 0.5756 & \itshape 0.5755 & \itshape 0.5757 & \itshape 0.9108 & \itshape 0.2185 \\
\quad Mistral, prompted (zero-shot) {\scriptsize (acc.)} &  & 0.2225 & 0.2173 & 0.2129 & 0.2176 & 0.5771 & 0.5795 & 0.5718 & 0.5761 & 0.9168 & 0.3674 \\
\midrule
\multicolumn{12}{l}{\itshape Reported for reference only; excluded from every statistic (see Section~\ref{subsec:res-setup})}\\
\itshape \quad Random ranking (floor) & 1 & \itshape 0.2181 & \itshape 0.2181 & \itshape 0.2181 & \itshape 0.2181 & \itshape 0.5820 & \itshape 0.5820 & \itshape 0.5820 & \itshape 0.5820 & \itshape --- & \itshape --- \\
\itshape \quad Llama, inv.\ $\lambda^\star$=0.05$^{\dagger}$ & 1 & \itshape 0.2228 & \itshape 0.2228 & \itshape 0.2228 & \itshape 0.2228 & \itshape 0.5938 & \itshape 0.5938 & \itshape 0.5938 & \itshape 0.5938 & \itshape --- & \itshape --- \\
\itshape \quad Mistral, English DPO$^{\dagger}$ & 1 & \itshape 0.2228 & \itshape 0.2228 & \itshape 0.2228 & \itshape 0.2228 & \itshape 0.5938 & \itshape 0.5938 & \itshape 0.5938 & \itshape 0.5938 & \itshape --- & \itshape --- \\
\itshape \quad Qwen, prompted (zero-shot)$^{\dagger}$ & 1 & \itshape 0.2228 & \itshape 0.2228 & \itshape 0.2228 & \itshape 0.2228 & \itshape 0.5938 & \itshape 0.5938 & \itshape 0.5938 & \itshape 0.5938 & \itshape --- & \itshape --- \\
\bottomrule
\end{tabular}
\caption{Condensed results on the held-out trilingual test set. Each training group shows its mean over $n$
configurations, then its most accurate (acc.) and, where different, its most cross-lingually stable
(stab.) member. Bold marks the best value per column among the $28$ valid configurations. Consistency
is Fleiss $\kappa$ over the language-conditioned verdicts and mean pairwise Kendall $\tau$ (EN--AR);
both are dashed for the reference rows, whose perfect agreement is an artefact of ranking
independently of the input (Section~\ref{subsec:res-setup}). Per-model results are in
Tables~\ref{tab:results-main} and~\ref{tab:results-consistency}.}
\label{tab:results-summary}
\end{table}


\paragraph{Every trained system beats the baselines, but the margins are narrow.}
The random ranking baseline reaches $0.2181$ Recall@5 and $0.5820$ macro-F1 in English, and the
zero-shot prompted baselines are barely distinguishable from it ($0.2234$ and $0.5751$ for the Llama
prompted baseline). This reproduces, on a larger and fully parallel evaluation set, the finding from
our shared-task system that zero-shot generative ranking adds little over chance in this setting. The
fine-tuned trace scorers improve substantially on macro-F1 (up to $0.6263$ in English) and modestly
on Recall@5 (up to $0.2710$).

\paragraph{Ranking headroom is nearly exhausted; verdict quality is not.}
The best Recall@5 of $0.2710$ corresponds to roughly $90\%$ of the achievable ceiling of $0.3005$. In
other words, the ranking component of the task is close to saturated for this candidate-trace set:
further gains cannot come from better ranking, only from candidate sets that contain more
gold-matching traces. Macro-F1, by contrast, ranges from $0.5715$ to $0.6320$ across configurations
and remains the component where modelling choices still matter. This asymmetry should be kept in mind
when reading the official combined score, which averages the two: it is increasingly dominated by the
verdict term.

\paragraph{Metric redundancy.}
Across the $28$ configurations, Recall@5, Precision@5, MRR@5, macro-F1, and the per-class F1 scores
for \textsf{True} and \textsf{False} are almost collinear, correlating at Pearson $r$ between $0.92$
and $0.99$. The single exception is F1 on the minority \textsf{Conflicting} class, which correlates
with the others at only $r=0.61$--$0.85$. Most of the headline metrics therefore restate the same
ordering of systems, and \textsf{Conflicting} is the class that actually discriminates between them.
We return to this class in Section~\ref{subsec:res-language}, where it turns out to carry the entire
cross-lingual effect.

\subsection{Cross-Lingual Consistency}
\label{subsec:res-consistency}

The right-hand block of Table~\ref{tab:results-summary} reports the consistency measures, computed
independently of accuracy (per-model values in Table~\ref{tab:results-consistency}). The central result of this section is that the two families of consistency measure behave
completely differently, and that only one of them discriminates between systems.


\begin{figure}[htbp]
  \centering
  \includesvg[width=\linewidth]{results-agreement}
  \caption{Verdict agreement versus ranking agreement, averaged by training type. Fleiss $\kappa$ over
  the three language-conditioned verdicts is high ($0.86$--$0.94$) for every group and separates them
  weakly. Mean pairwise Kendall $\tau$ between the language-conditioned trace rankings separates the
  same groups by almost an order of magnitude.}
  \label{fig:results-agreement}
\end{figure}

\paragraph{Verdicts agree across languages; rankings do not.}
As Figure~\ref{fig:results-agreement} shows, Fleiss $\kappa$ over the language-conditioned verdicts is
uniformly high and occupies a narrow band: $0.8596$ for the earlier trace scorers, $0.9028$ for
standard fine-tuning, and $0.9443$ for invariance-trained models. Judged on verdict agreement alone,
every system looks strongly language-invariant. Kendall $\tau$ between the language-conditioned trace
\emph{rankings} tells a very different story: on the English--Arabic pair it spans $0.012$ to $0.644$
across individual configurations and $0.0578$ to $0.5667$ across groups, and on the English--Spanish
pair the weakest configuration is very slightly negative ($-0.008$), i.e.\ no better than an unrelated
ordering.

The interpretation is that standard multilingual fine-tuning produces models that reach the same
verdict in all three languages while justifying it with almost entirely different evidence: the
top-ranked traces in English and Arabic are close to uncorrelated ($\tau \approx 0.13$). Because the
final verdict is a majority vote over the top five traces, a model can arrive at a stable label
through unstable and largely disjoint evidence. Verdict-level agreement statistics therefore
substantially overstate cross-lingual robustness, and ranking-level agreement is the measure that
exposes the difference. This is a methodological finding in its own right: had we followed the common
practice of reporting only $\kappa$ over final answers, we would have concluded that all our systems
were already language-invariant.

\subsection{Does Invariance Training Improve Ranking? (RQ1)}
\label{subsec:res-invariance}

Figure~\ref{fig:results-tradeoff} places every configuration in the plane defined by mean macro-F1
(accuracy) and Kendall $\tau$ between English and Arabic rankings (stability). The picture is
unusually clean.

\begin{figure}[htbp]
  \centering
  \includesvg[width=\linewidth]{results-tradeoff}
  \caption{Each point is one configuration ($n=28$). The vertical axis is the mean pairwise Kendall
  $\tau$ between the English and Arabic trace rankings; the horizontal axis is macro-F1 averaged over
  the three languages. The invariance-trained population is fully separated from every other
  configuration on stability, while the accuracy ranges overlap.}
  \label{fig:results-tradeoff}
\end{figure}


\paragraph{The stability effect is large and fully separated.}
Table~\ref{tab:results-summary} summarizes the group-level picture.
The $13$ invariance-trained configurations reach a mean Kendall $\tau$ of $0.5667$ between English and
Arabic, against $0.1308$ for standard fine-tuning, $0.0578$ for the earlier trace scorers, and
$0.2185$ for the baselines. The two populations do not overlap at all: the \emph{worst}
invariance-trained configuration ($\tau = 0.5089$) exceeds the \emph{best} non-invariance
configuration ($\tau = 0.3674$) by a wide margin. The effect size is Cohen's $d = 6.09$
(Welch $t = 16.71$, $p < 10^{-13}$; Mann--Whitney $p < 10^{-5}$). Verdict agreement moves in the same
direction, from Fleiss $\kappa = 0.8924$ to $0.9443$ ($d = 2.20$, $p < 10^{-4}$). Complete separation
of two populations is a stronger statement than any $p$-value computed from $28$ points, and it is the
claim we rest on.

\paragraph{The accuracy effect is positive but not the argument.}
Mean macro-F1 is also higher for invariance-trained models, $0.6148$ against $0.5977$
($d = 1.61$, $p < 0.001$). However, unlike stability, the accuracy \emph{ranges overlap}: the best
non-invariance configuration (\emph{Llama, standard}, $0.6227$) outperforms the best
invariance-trained one (\emph{Qwen, inv.\ $\lambda^\star$=0.1}, $0.6182$). If invariance training were
judged on English macro-F1 alone it would read as a null result, with the best English score
belonging to \emph{Llama, reg.} ($0.6263$).

\paragraph{On the official ranking metric, the two are tied.}
Recall@5 deserves separate comment, because it is the official ranking measure and because it does
\emph{not} follow macro-F1. Averaged over the three languages, standard fine-tuning attains a
marginally \emph{higher} Recall@5 than invariance training ($0.2679$ against $0.2656$), and the single
best Recall@5 of the study belongs to a standard fine-tune (\emph{Llama, reg.}, $0.2710$ in English).
Per language the two groups are within $0.007$ of each other everywhere
(Figure~\ref{fig:results-language}, right panel). Given the ceiling of $0.3005$ and the near-collinearity
of the ranking metrics noted in Section~\ref{subsec:res-task}, the fair reading is that the two
approaches are tied on Recall@5: invariance training neither buys nor costs anything on the amount of
relevant evidence retrieved into the top five.

The honest answer to RQ1 is therefore conditional, and it is worth stating precisely. Optimizing for
language invariance does \emph{not} produce a better ranker in the sense of a higher per-language
Recall@5; on that metric it is level with standard fine-tuning, and slightly behind on average. What
it produces is a ranker whose ordering of reasoning traces is stable under translation---a roughly
fourfold improvement in cross-lingual ranking agreement---together with better verdict quality
(macro-F1) and a markedly smaller loss into Arabic. Since the property being optimized is
cross-lingual consistency rather than per-language retrieval, this is the outcome the objective was
designed to produce; but it means the contribution should be claimed as robustness, not as
state-of-the-art ranking performance.

\subsection{Where the Cross-Lingual Differences Live}
\label{subsec:res-language}

\begin{figure}[htbp]
  \centering
  \includesvg[width=\linewidth]{results-language}
  \caption{Both official metrics by language, averaged by training type. \emph{Top:} macro-F1.
  Spanish is on par with English for every group, Arabic is consistently the weakest, and the Arabic
  gap is smallest for invariance-trained models. \emph{Bottom:} Recall@5 against its attainable
  ceiling of $0.3005$ (dashed). The two fine-tuned groups are tied on Recall@5 and sit close to the
  ceiling, so the invariance effect visible in the upper panel does not appear in the lower one.}
  \label{fig:results-language}
\end{figure}

Averaged over the $28$ configurations, Spanish macro-F1 ($0.6118$) is marginally \emph{above} English
($0.6099$), while Arabic ($0.5951$) sits $0.0147$ below English, a relative drop of about $2.4\%$.
Recall@5 follows the same pattern ($0.2579$ / $0.2620$ / $0.2548$). The multilingual difficulty in
this task is thus almost entirely an Arabic effect, consistent with the expectation from
Section~\ref{subsec:bg-numeracy} that a different numeral script and lower pre-training representation
would make Arabic the hard case~\cite{reddy2026effectscriptsformatsllm,Qi_2023}.

\begin{figure}[htbp]
  \centering
  \includesvg[width=\linewidth]{results-conflicting}
  \caption{F1 on the minority \textsf{Conflicting} class by language. The Arabic degradation seen in
  Figure~\ref{fig:results-language} is concentrated almost entirely in this class: standard fine-tunes
  fall from $0.270$ in English to $0.178$ in Arabic, while invariance-trained models hold at $0.251$.}
  \label{fig:results-conflicting}
\end{figure}

Figure~\ref{fig:results-conflicting} localizes the effect. The aggregate Arabic drop is not spread
evenly over the label space. F1 on \textsf{False} and \textsf{True} is almost flat across languages,
losing only $0.0026$ and $0.0027$ respectively between English and Arabic, whereas F1 on the minority
\textsf{Conflicting} class falls from $0.2625$ (English) and $0.2671$ (Spanish) to $0.2236$
(Arabic)---a drop roughly fourteen times larger than for either majority class. Broken down by training type the contrast is sharper still.
Averaging the per-configuration \textsf{Conflicting} F1 within each group and language, standard
fine-tunes drop from $0.2701$ in English to $0.1781$ in Arabic---a relative loss of $34.1\%$ of the
English value---whereas invariance-trained models fall only from $0.2698$ to $0.2513$, a relative loss
of $6.8\%$. Unlike the aggregate Arabic gap
discussed below, this difference is statistically clear: Welch $p = 0.0001$, Mann--Whitney
$p = 0.0003$, with only two of the nine standard fine-tunes reaching into the range spanned by the
thirteen invariance-trained models. \textsf{Conflicting} is both the rarest class ($427$ of $4{,}233$
claims, $10.1\%$) and the one that requires genuinely weighing contradictory quantitative evidence
rather than matching a single number, so it is the class where unstable evidence selection is most
damaging. The class-level view thus explains the aggregate: invariance training helps Arabic
specifically by protecting the hardest and rarest decision.

\paragraph{A caution on the aggregate Arabic gap.}
The mean English-to-Arabic macro-F1 drop is $-0.0098$ for invariance-trained models against $-0.0311$
for standard fine-tuning, a difference of roughly threefold in the expected direction. We note,
however, that this particular comparison does \emph{not} reach statistical significance
(Welch $p = 0.089$; Mann--Whitney $p = 0.231$), because the per-configuration variance in the gap is
large relative to $n$. The direction is consistent across families and settings and is corroborated by
the class-level analysis above, but it should be reported as a consistent trend rather than an
established effect. This contrasts with the ranking-stability result, which is unambiguous.

\subsection{Ablations}
\label{subsec:res-ablations}

\paragraph{The invariance weight.}
Figure~\ref{fig:results-lambda} sweeps the target invariance weight $\lambda^\star$ over
$\{0.1, 0.2, 0.3\}$ for both backbones. Cross-lingual stability increases monotonically with
$\lambda^\star$ for both families---Llama from $0.5711$ to $0.6440$ and Qwen from $0.5820$ to
$0.6426$---which is the expected behaviour of a regularizer whose strength is being increased. Mean
macro-F1 does not follow: it peaks at $\lambda^\star = 0.2$ for Llama ($0.6180$, against $0.6146$ and
$0.6116$ at $0.1$ and $0.3$) and declines slightly but monotonically for Qwen ($0.6182 \to 0.6173 \to
0.6152$). Recall@5 behaves like macro-F1 rather than like $\tau$, peaking at $\lambda^\star = 0.2$ for
both backbones. The sweep therefore
exposes a mild trade-off rather than a free lunch, and $\lambda^\star = 0.2$ is the setting that
balances the two. It should be stressed that the entire swept range remains far above the
non-invariance population on stability, so the choice of $\lambda^\star$ within this range is much less
consequential than the decision to use the objective at all.

\begin{figure}[htbp]
  \centering
  \includesvg[width=\linewidth]{results-lambda}
  \caption{Effect of the target invariance weight $\lambda^\star$. Cross-lingual ranking stability
  rises monotonically with $\lambda^\star$ for both backbones, whereas both task metrics---macro-F1
  and Recall@5---peak at $\lambda^\star = 0.2$. The objective therefore trades a little task quality
  for consistency beyond that point, making $\lambda^\star = 0.2$ the balanced choice.}
  \label{fig:results-lambda}
\end{figure}

\paragraph{Numeric normalization does not combine with invariance training.}
Plain-text numeric normalization was carried over from the shared-task system as a complementary lever
for the numerical dimension of the task (Section~\ref{subsec:method-numeric}). Within the
invariance-trained group it does not help. Configurations combining the two reach $\tau = 0.5121$ and
mean macro-F1 $0.6133$, against $\tau = 0.5910$ and $0.6154$ for invariance training alone: adding
numeric normalization costs $0.079$ Kendall $\tau$ and buys no accuracy, and the same pattern holds in
Arabic ($0.6028$ against $0.6111$). Outside the invariance group the picture is different---among
standard fine-tunes, numeric normalization is mildly positive in English and Spanish---so the two
interventions appear to be partially redundant rather than individually useless. A plausible reading is
that both mechanisms address the same underlying problem: numeric normalization canonicalizes numerals
so that they are represented identically across scripts, which is one of the specific ways the
invariance objective already forces scores to agree. Once the score-level constraint is imposed,
restating canonical values in the input adds input length without adding signal. On this evidence, the
two should not be combined.

\paragraph{Backbone families.}
Pooling all settings, Qwen3-8B and Llama-3.1-8B are level on accuracy (mean macro-F1 $0.6128$ against
$0.6083$) but Qwen is noticeably more stable cross-lingually ($\tau = 0.4505$ against $0.3161$;
Fleiss $\kappa = 0.9353$ against $0.9195$). Mistral-7B trails on both dimensions ($0.5843$ and
$0.1197$), confirming the earlier decision to retain it as a reference point rather than extend it with
the invariance objective. The strongest individual results are split: \emph{Llama, reg.} attains the
best English macro-F1 ($0.6263$) and Recall@5 ($0.2710$), \emph{Llama, inv.\ $\lambda^\star$=0.2} the
best Arabic macro-F1 ($0.6161$), and \emph{Llama, inv.\ $\lambda^\star$=0.3} the best ranking stability
($\tau = 0.6440$). That no single configuration wins on every axis is itself consistent with the
accuracy/stability distinction drawn in Section~\ref{subsec:res-invariance}.

\subsection{The Relationship Between Invariance and Reasoning Quality (RQ2)}
\label{subsec:res-relationship}

RQ2 asks how language invariance relates to reasoning quality. Because the two were measured
independently, the relationship can be examined empirically. Throughout this section we quantify
reasoning quality by \emph{mean macro-F1}, the claim-verdict macro-F1 averaged over the three
languages, and we report the ranking metric separately rather than folding both into a single notion
of ``accuracy'', because---as the numbers below show---the two behave differently.

Over the $28$ configurations, mean macro-F1 correlates with Kendall $\tau$ at $r = 0.60$ and with
Fleiss $\kappa$ at $r = 0.57$. The corresponding correlations for the ranking metric are weaker:
mean Recall@5 correlates with $\tau$ at only $r = 0.39$ (English Recall@5 alone, $r = 0.49$), while
still correlating with $\kappa$ at $r = 0.57$. The association between quality and consistency is
therefore positive and moderate, as the thesis premise anticipated, but it holds mainly for verdict
quality; how much relevant evidence a model retrieves into its top five says comparatively little
about whether it retrieves the same evidence in another language.

The association is, however, far from deterministic, and Figure~\ref{fig:results-tradeoff} makes the
residual visible. At essentially identical mean macro-F1 (around $0.61$--$0.62$) the configurations
span $\tau$ from $0.08$ to $0.64$---a range of more than $0.5$ at fixed verdict quality. Mean macro-F1
consequently predicts stability only weakly, and the two properties are separable in practice as well
as in principle. This is the empirical confirmation, in a ranking setting, of the decoupling that the
consistency literature established for factual
prediction~\cite{elazar2021measuringimprovingconsistencypretrained,Qi_2023}: one cannot infer that a
more accurate multilingual model is a more consistent one, and improving task quality is not a route
to consistency.

A second relationship is worth recording. Fleiss $\kappa$ and Kendall $\tau$ correlate at $r = 0.83$
across configurations, so the two consistency families are related---but, as
Section~\ref{subsec:res-consistency} showed, $\kappa$ is compressed into a narrow high band while
$\tau$ is not. A high correlation between two measures does not make them interchangeable when one has
far less dynamic range than the other.

\subsection{Threats to Validity}
\label{subsec:res-threats}

\paragraph{Context-budget confound.}
The most important qualification concerns training context. The invariance-trained models were trained
with a reduced context budget, because the aligned trilingual batching required for the objective
multiplies the sequence count per optimization step and the available compute did not permit the full
long-context configuration. Several of the standard fine-tunes, by contrast, used the full budget.
Training condition and context budget are therefore partially entangled, and the comparison is not a
controlled ablation.

Two considerations limit the damage. First, the direction of the confound is conservative: the
invariance arm is the handicapped one, so its wins are lower bounds. Second, the size of the context
effect can be estimated from configurations that were run at both budgets. Moving the same
architectures from the reduced to the full budget changes English macro-F1 by $+0.0052$, Spanish by
$+0.0043$, and Arabic by $-0.0053$, for a net change of $+0.0014$ on the three-language mean---smaller
than the invariance effect and not even consistent in sign across languages. Restricting the
comparison to the matched reduced-budget condition, the invariance-trained models still exceed the
short-context fine-tunes by $+0.0083$ mean macro-F1, $+0.0197$ Arabic macro-F1, and a factor of $4.85$
on Kendall $\tau$. No invariance-trained configuration was run at the full context budget; that single
run would resolve the confound and is the most valuable next experiment. We therefore state the claim
in the guarded form that the evidence supports:
\emph{invariance training at a reduced context budget matches or exceeds full-context fine-tuning on
mean macro-F1, and dominates it on cross-lingual stability}. The stability result is
unaffected by the confound, since no plausible context effect of the magnitude estimated here could
produce a fourfold change in $\tau$.

\paragraph{Small comparison arms.}
The matched-budget comparison rests on only two short-context fine-tunes (one Llama, one Qwen), and the
full-budget arm has seven configurations. The group means are correspondingly imprecise, and the
Spanish result is an explicit exception to the general pattern: \emph{Llama, short context} reaches
$0.6220$ Spanish macro-F1, above all $13$ invariance-trained configurations. We do not claim a Spanish
improvement. The defensible summary is parity on English and Spanish, a gain on Arabic, and a large gain
on cross-lingual stability.

\paragraph{Behavioural, not representational, evidence.}
Finally, all invariance measures reported here are behavioural: they compare outputs across languages.
As noted in Section~\ref{subsec:bg-multilingual}, agreement of outputs does not establish that the model
represents the underlying knowledge in a shared way~\cite{ifergan2024beneathsurfaceconsistencyexploring}.
Our results show that the objective changes what the model \emph{does} across languages; they do not
establish what it has changed internally.

\subsection{Summary of Findings}
\label{subsec:res-summary}

\begin{description}
\item[RQ1 --- Does optimizing for language-invariant reasoning improve ranking performance?]
Partly, and in a specific sense. On the official ranking metric the objective changes nothing: mean
Recall@5 is $0.2656$ for invariance training against $0.2679$ for standard fine-tuning, and the best
Recall@5 in the study ($0.2710$) belongs to a standard fine-tune. What it does produce is a fourfold,
fully separated improvement in the cross-lingual stability of the ranking
($\tau = 0.567$ against $0.131$, $d = 6.09$), together with a small but consistent gain in verdict
quality ($0.6148$ against $0.5977$ mean macro-F1). Ranking quality measured per language is essentially unchanged; ranking
\emph{consistency} across languages is transformed.

\item[RQ2 --- What is the relationship between language invariance and reasoning quality?]
Positive but moderate and clearly non-deterministic. Mean macro-F1---verdict quality averaged over the
three languages---correlates with ranking stability (Kendall $\tau$) at $r=0.60$ and with verdict
agreement (Fleiss $\kappa$) at $r=0.57$; the ranking metric is a weaker predictor still, with mean
Recall@5 correlating with $\tau$ at only $r=0.39$. Yet at fixed mean macro-F1 the observed stability
spans a range of more than $0.5$ in $\tau$. Consistency is not a by-product of task quality, and the
two must be measured separately---a conclusion reinforced by the random baseline, which is perfectly
invariant and entirely uninformative.

\item[RQ3 --- Can language-invariant training objectives lead to more robust multilingual reasoning
models?] Yes, with the qualifications above. The gains are concentrated where robustness is most needed:
in Arabic, the different-script and least well-represented of the three languages, and within Arabic in
the minority \textsf{Conflicting} class, where standard fine-tuning degrades several times more than
invariance training. The effect holds across both Llama and Qwen backbones, indicating a property of
the method rather than of a single model family, and it is obtained under a reduced training-context
budget.
\end{description}

A final observation cuts across all three questions. Standard multilingual fine-tuning already produces
models that agree on the verdict in every language; what it does not produce is models that agree on
\emph{why}. The contribution of the invariance objective is to align the evidence, not the label---and
because that difference is invisible to verdict-level agreement statistics, it would have gone unmeasured
under the evaluation protocol most commonly used in this literature.
